Esta interpretación la redactó Claude a partir de los registros de las corridas; no es del autor.

Luna no aceptó ninguno de los 12 episodios de los cuatro objetivos (48 intentos). No es que el juez sea inalcanzable: para los dos objetivos con $n$ fijo, la demostración generada desde el certificado exacto del PEP es \textsc{Accepted} en unos 5\,s (Observación~2). Lo que falta es escribirla. En \texttt{gd\_n1}, 4 de los 12 intentos se perdieron por el analizador de respuestas anterior (nota de la carpeta); los otros 8 son demostraciones rechazadas. En \texttt{gd\_const\_textbook}, el único \textsc{Statement\_Drift} es una comilla sin cerrar en el bloque de lemas: el filtro trata todo texto que no puede leer como un intento de salir del hueco, pero aquí se trata de un error de sintaxis, no de un intento de cambiar el enunciado. Los tres \textsc{Timeout} de \texttt{gd\_const\_dt} son compilaciones que pasaron de 300\,s.

Los errores repiten el patrón de la bisección: métodos de Lean y Mathlib que Isabelle no tiene (\texttt{nlinarith} 6 veces, \texttt{ring} 3, \texttt{positivity} 1) y nombres de hechos inexistentes, casi todos con forma de nombres de Mathlib (\texttt{sq\_nonneg}, \texttt{inner\_self\_eq\_norm}, \texttt{inner\_le\_norm}); el hecho que Luna buscaba en Isabelle se llama \texttt{power2\_norm\_eq\_inner}.

Hay un resultado positivo. En la cota de Drori y Teboulle, Luna demostró e Isabelle verificó dos lemas propios: el desdoblamiento de la iteración,
\[ \mathrm{gd}\ g\ (\mathrm{replicate}\ (n+1)\ 1)\ x = \mathrm{gd}\ g\ (\mathrm{replicate}\ n\ 1)\ (x - g(x)), \] y la co-coercividad, $\langle g(x) - g(y), x - y\rangle \ge \|g(x) - g(y)\|^2$, obtenida sumando las dos instancias de \texttt{interp} con una demostración Isar estructurada. La co-coercividad es un paso estándar de las demostraciones de libro de las cotas del descenso de gradiente, y es el primer lema verificado en las corridas de stepsize (en la bisección, los lemas verificados de Luna eran hechos auxiliares sobre logaritmos y techos). No bastó: la cota exacta para todo $n$ necesita la combinación completa de desigualdades con los multiplicadores de Drori y Teboulle.

Aquí el juez cuesta unos 5\,s por intento y el modelo unos 3\,min: a diferencia de la bisección, el cuello de botella es el proponente.
